% Citation map verified 2026-07-24 (all arXiv IDs resolve); keys in iclr2026_conference.bib.

\paragraph{RL fine-tuning of generative robot policies.}
Policy-gradient fine-tuning of vision-language-action models with sparse task
reward is now standard for autoregressive token policies
\citep{simplevlarl}, but transfers poorly to continuous generative heads:
likelihood ratios are intractable for diffusion and flow policies, and
workarounds either reformulate the denoising chain as an MDP
\citep{ren2024dppo,black2023ddpo} --- a construction with no analog for
one-step generators --- or avoid likelihoods entirely
\citep{park2025fql,reinflow2025,parl2024}. FPO
\citep{mcallister2025flowmatchingpolicygradients} is the strongest form of
the likelihood-free ratio: it replaces $\pi_\theta/\pi_{\theta_{\text{old}}}$
with the exponentiated difference of per-sample conditional flow-matching
losses, which the CFM--ELBO relationship makes a valid proxy, and it does
so without perturbing the deployed policy. It is therefore the right
policy-gradient method for a flow or diffusion policy --- but the
surrogate is defined by a denoising loss over a flow time $\tau$, and a
one-step drifting model has neither: its training objective is
field-target regression, with no ELBO to bound. FPO also fine-tunes (or
trains from scratch) the full policy on continuous-control suites, where
we update a residual on a frozen base for manipulation. The two are
complementary --- FPO makes the \emph{ratio} tractable, we make the
\emph{critic} safe to follow --- and we regard evaluating FPO on a
flow-matching base against our update as the clearest missing comparison
in this paper. Closest to our setting, DICE-RL
\citep{dicerl2026} fine-tunes a frozen one-step distillation of a
flow-matching policy by training a residual head with a backpropagated
critic and a filtered BC anchor; it is our primary baseline and the
\emph{control} for our actor-update ablation, since we keep its critic,
residual parameterization, and data pipeline unchanged.

\paragraph{Where this update sits.}
RL for diffusion and flow policies divides into four families
\citep{dprlsurvey2026}: backpropagating through the denoising chain
(Diffusion-QL, and DICE-RL's residual actor); ratio-based methods that make
a likelihood tractable (DPPO, FPO); $Q$-weighted regression; and
\emph{action-gradient} methods, which improve actions with $\nabla_a Q$ and
then regress the policy onto the improved actions. \method{} belongs to the
last family, and we want to be exact about what is and is not new in it.
The primitive of displacing actions along $\nabla_a Q$ and regressing onto
the result is \emph{not} ours: DIPO \citep{dipo2024} introduced it, updating
replay-buffer actions by gradient ascent and fitting a diffusion policy to
them, and QSM \citep{psenka2024qsm} reaches the same signal differently, by
matching the policy's score to $\nabla_a Q$. What the family lacks --- and
what the survey notes none of its members supply --- is any bound on how far
that displacement may carry the policy, or any restoring force toward a
behavior policy. That gap is also its documented weakness: action-gradient
methods are reported to degrade when the critic's gradient is biased, and to
drift away from the behavior distribution.

Our contribution is precisely that missing piece, and our ablations measure
it directly. The displacement is clipped pointwise in action space, a
dead-zone anchor restores toward a frozen base, and the update is applied to
a residual rather than by overwriting stored transitions (DIPO's buffer
rewrite makes the stored action inconsistent with the transition that
produced it). Two measurements tie our system to theirs. Removing the
base-restoring anchor reproduces the family's failure inside our own code
exactly --- success $0.000$, whether or not the pointwise clip is kept
(\cref{tab:constraint}) --- so we can switch their failure on and off. And
running DIPO, QSM and DQL from their authors' implementations on our setting,
at $125\times$ our interaction budget, none improves its base and two reach
exactly $0.000$ (\cref{tab:baselines}), which is the same failure observed
from the outside. We are careful about the direction of this claim: these
methods were developed for dense-reward online control, and the adaptation to
sparse-reward manipulation fine-tuning is the one published by
\citet{ren2024dppo}, not one we devised. Their difficulty here is evidence
about the regime, not a verdict on the methods. Test-time guidance
\citep{qgf2026} applies the same $\nabla_a Q$ signal without training at all,
and is complementary.

\paragraph{Critic-as-compass versus critic-as-judge.}
That a learned $Q$-function's action-gradient cannot be trusted is a
recurring diagnosis: overestimation and narrow spurious peaks in TD3
\citep{fujimoto2018td3}, extrapolation error in BCQ
\citep{fujimoto2019bcq}, and spurious local optima that make
sample-and-select beat gradient ascent \citep{savo2025}. One response
treats the critic only as a \emph{judge}: weighted regression and
sample-based improvement (AWR, CRR, MPO, IDQL, QT-Opt)
\citep{peng2019awr,wang2020crr,abdolmaleki2018mpo,hansenestruch2023idql,kalashnikov2018qtopt}
tilt the policy toward $\pi \cdot \exp(Q/\tau)$ using only $Q$-values.
Our method spans this spectrum rather than picking a side: the update
\emph{form} is judge-safe (bounded transport toward targets), while the
reward field inside it can still be the analytic gradient where it is
trustworthy --- and our ablations map exactly when each choice wins. The
judge-only methods are also the ones that survive our benchmark: of the
published fine-tuners we ran, the one that only \emph{weights} actions rather
than displacing them is the only one to improve its base
(\cref{tab:baselines}). It pays for that safety by discarding the critic's
most informative output, which is what our bounded update is designed to
keep.

\paragraph{Particle transport as a policy update.}
Regressing a sampler onto displaced particles is the amortized-SVGD idea
of Soft Q-Learning \citep{haarnoja2017sql}, with kernel repulsion
preserving diversity as in SVPG \citep{liu2017svpg}; drifting models
\citep{deng2026drifting} make the same operator the \emph{pretraining}
objective of a one-step generator. Drifting Field Policy
\citep{dfp2026} is the closest prior method: it also performs RL on a
drifting backbone by drifting fresh policy samples toward
critic-ranked top-$K$ candidates. \method{} differs in three ways: we
update a \emph{residual on a frozen base} rather than the full policy,
we transport with an explicit pointwise trust region and a restoring
anchor field (stability contract our toy shows is load-bearing),
and our reward field is a dial --- including the analytic
$\nabla_a Q$ source that dominates when the critic generalizes ---
rather than a fixed ranking surrogate.

\paragraph{Residual RL and staying on-manifold.}
Learning a residual on a frozen base is a long-standing recipe for
sample-efficient robot RL; DAWN \citep{dawn2026} identifies the value
cold-start and scale pathologies of residual critics that we also
observe. The off-manifold drift our toy exhibits is the mechanism
behind over-optimization and mode collapse in RL-finetuned generators
\citep{black2023ddpo,fan2023dpok,nakamoto2023calql}, and forgetting in
RL-finetuned language models is likewise governed by divergence from
the base \citep{rlrazor2025} --- which cannot even be measured without
likelihoods, reinforcing update rules whose step size is bounded in
\emph{action} space by construction.

In sum, prior work either trusts the critic's gradient and pays for it
off-manifold, or discards it and pays on hard tasks; \method{} is, to our
knowledge, the first fine-tuner for one-step generative policies whose
update \emph{form} is safe under an arbitrary critic while still admitting
the full first-order signal when it is available.
